\documentclass[a4paper,12pt]{article}
\usepackage[T2A]{fontenc}
\usepackage[utf8]{inputenc}
\usepackage[russian]{babel}
\usepackage{amsmath,amssymb,icomma}
\usepackage[a4paper,margin=20mm,headheight=25pt,headsep=8mm,footskip=12mm]{geometry}
\usepackage{microtype,xcolor,tikz,pgfplots,booktabs,tabularx,enumitem,fancyhdr,setspace}
\usetikzlibrary{arrows.meta,positioning,shapes.geometric,decorations.pathmorphing}
\pgfplotsset{compat=1.18}
\definecolor{accent}{HTML}{008080}
\definecolor{darkaccent}{HTML}{004D4D}
\definecolor{textgray}{HTML}{333333}
\color{textgray}
\setstretch{1.35}
\setlength{\parindent}{0pt}
\setlength{\parskip}{5pt}
\setlist[itemize]{leftmargin=*,itemsep=1pt,topsep=3pt}
\setlist[enumerate]{leftmargin=*,itemsep=4pt,topsep=3pt}
\pagestyle{fancy}
\fancyhf{}
\fancyhead[L]{\small Механика: кинематика и динамика}
\fancyhead[R]{\small Володя Хачатурян}
\fancyfoot[C]{\small\thepage}
\renewcommand{\headrule}{\hbox to\headwidth{\color{accent}\leaders\hrule height 0.5pt\hfill}}
\newcommand{\topic}[1]{\par\vspace{2pt}{\Large\color{darkaccent}#1}\par\vspace{4pt}{\color{accent}\hrule height 0.4pt}\vspace{8pt}}
\newcommand{\sub}[1]{\par\vspace{5pt}{\color{darkaccent}\textbf{#1}}\par}
\newcommand{\unit}[1]{\,\text{#1}}
\tikzset{vec/.style={-{Stealth[length=2.5mm]},very thick,accent},axis/.style={-{Stealth[length=2mm]},thin},spring/.style={decorate,decoration={coil,aspect=.45,segment length=4pt,amplitude=3pt},thick,accent},action/.style={draw=accent,rounded corners=2pt,align=center,text width=10.5cm,minimum height=12mm,inner sep=5pt},terminal/.style={draw=accent,ellipse,align=center,minimum width=4cm,minimum height=9mm},flowarrow/.style={-{Stealth[length=2.5mm]},draw=darkaccent,thick}}
\begin{document}
\begin{titlepage}
\centering
\vspace*{26mm}
{\Large\color{darkaccent} ФИЗИКА\par}
\vspace{14mm}
{\huge Чек-лист\par}
\vspace{10mm}
{\LARGE\color{accent} Механика:\\ кинематика и динамика\par}
\vspace{14mm}
{\large Решение типовых задач и анализ движения\par}
\vfill
{\large Володя Хачатурян\par}
\vspace{9mm}
{25 ноября 2026 г.\par}
\vspace{17mm}
{\small Лёвин Артём Александрович\\ эксклюзивно для Володи Хачатуряна\par}
\vspace{17mm}
\begin{tikzpicture}
\draw[accent,line width=.65pt] (0,0) .. controls (2,0.8) and (4,-0.8) .. (6,0) .. controls (8,0.8) and (10,-0.8) .. (12,0);
\end{tikzpicture}
\vspace*{5mm}
\end{titlepage}
\setcounter{page}{1}

\topic{1. Равноускоренное прямолинейное движение}
\textbf{Главная идея:} сначала опишите движение и выберите координатную ось, затем запишите уравнения в проекциях. При постоянном ускорении $a_x$:
\[
 v_x=v_{0x}+a_xt,\qquad x=x_0+v_{0x}t+\frac{a_xt^2}{2}.
\]
Здесь $x_0$ --- начальная координата, $v_{0x}$ --- проекция начальной скорости, $t$ --- время. При $a_x=0$ скорость постоянна.
\sub{Как извлечь данные из уравнения координаты}
Если дано $x=A+Bt+Ct^2$, сравните коэффициенты:
\[
 x_0=A,\qquad v_{0x}=B,\qquad a_x=2C.
\]
Пример: $x=-11+3t+t^2$ (величины в СИ). Коэффициент при $t^2$ равен единице, поэтому $a_x=2\unit{м/с$^2$}$, $v_{0x}=3\unit{м/с}$, $x_0=-11\unit{м}$.
\[
 v_x=3+2t,\qquad v_x(5)=3+2\cdot5=13\unit{м/с}.
\]
\sub{Движение из состояния покоя}
Если шарик скатывается с постоянным ускорением, $v_{0x}=0$, ось направлена вдоль движения. За первую секунду он прошёл 5 см. Переводим в СИ: $5\unit{см}=0{,}05\unit{м}$.
\[
 s_x=\frac{a_xt^2}{2}\quad\Longrightarrow\quad a_x=\frac{2s_x}{t^2}
 =\frac{2\cdot0{,}05}{1^2}=0{,}1\unit{м/с$^2$}.
\]
\sub{Проверяйте}
\begin{itemize}
\item Начальная скорость может быть нулевой, а начальная координата --- произвольной.
\item Коэффициент при $t^2$ равен $a_x/2$.
\item Знаки проекций определяются направлением выбранной оси.
\end{itemize}
\clearpage

\topic{2. Перемещение, путь и графики}
Перемещение --- вектор от начального положения к конечному. При движении по одной прямой его проекция и модуль равны
\[
 s_x=x_{\text{кон}}-x_{\text{нач}},\qquad |\vec s|=|s_x|.
\]
Путь $S$ --- суммарная длина пройденных участков. При смене направления движения он может превышать модуль перемещения.
\sub{Пример 1. Возвращение в начальную точку}
\begin{center}
\begin{tikzpicture}
\begin{axis}[width=10cm,height=4.6cm,axis lines=left,xmin=0,xmax=6.6,ymin=-25,ymax=25,xtick={0,1,4,6},ytick={-20,0,20},xlabel={$t,\ \text{с}$},ylabel={$x,\ \text{м}$},grid=major,grid style={gray!18},tick label style={font=\small},label style={font=\small}]
\addplot[accent,very thick,mark=*] coordinates {(0,-20) (1,20) (4,20) (6,-20)};
\end{axis}
\end{tikzpicture}
\end{center}
По графику: $x(0)=-20\unit{м}$ и $x(6)=-20\unit{м}$. Тогда $|\vec s|=|-20-(-20)|=0$, хотя пройденный путь равен $40+40=80\unit{м}$.
\sub{Три формулы для проекции перемещения}
При постоянном ускорении:
\[
 s_x=v_{0x}t+\frac{a_xt^2}{2},\qquad
 s_x=\frac{v_x^2-v_{0x}^2}{2a_x}\ (a_x\ne0),
\]
\[
 s_x=\frac{v_{0x}+v_x}{2}\,t.
\]
При движении в одном направлении $S=|s_x|$. При смене направления путь считают по отдельным участкам.
\sub{Пример 2. Тормозной путь}
Скорость равномерно уменьшается от 15 м/с до нуля за 20 с. Выбираем третью формулу: ускорение вычислять отдельно не требуется.
\[
 S=\frac{15+0}{2}\cdot20=150\unit{м}.
\]
На графике $v_x(t)$ проекция перемещения равна площади с учётом знака; при таком торможении это площадь треугольника.
\clearpage

\topic{3. Вертикальное движение}
При пренебрежимо малом сопротивлении воздуха ускорение свободного падения направлено вниз. Его модуль $g$ положителен. Если ось $Oy$ направлена вверх, $a_y=-g$:
\[
 v_y=v_{0y}-gt,\qquad y=y_0+v_{0y}t-\frac{gt^2}{2}.
\]
\sub{Бросок вертикально вверх}
\begin{center}
\begin{tikzpicture}[scale=.88]
\draw[thick] (-1.8,0)--(3.4,0);
\foreach \x in {-1.7,-1.3,...,3.1} \draw[gray] (\x,0)--++(-.23,-.23);
\draw[axis] (2.8,0)--(2.8,3.3) node[above] {$y$};
\draw[dashed,gray] (0,2.75)--(2.8,2.75);
\fill[accent] (0,.12) circle (2.2pt); \fill[accent] (0,2.75) circle (2.2pt);
\draw[vec] (0,.25)--(0,1.45) node[right] {$\vec v_0$};
\draw[vec] (1.5,2.5)--(1.5,1.3) node[right] {$\vec g$};
\draw[<->] (-.8,.04)--(-.8,2.75) node[midway,left] {$h_{\max}$};
\node[above] at (0,2.75) {$v_y=0$};
\end{tikzpicture}
\end{center}
В верхней точке скорость вертикального броска равна нулю, ускорение сохраняется: $a_y=-g$.
\[
 t_{\max}=\frac{v_0}{g},\qquad h_{\max}=\frac{v_0^2}{2g}.
\]
Пример: тело бросили с земли со скоростью 4 м/с, $g=10\unit{м/с$^2$}$. Сначала находим время подъёма из условия $v_y=0$:
\[
 4-10t=0\ \Rightarrow\ t=0{,}4\unit{с},\qquad
 h_{\max}=4\cdot0{,}4-5\cdot(0{,}4)^2=0{,}8\unit{м}.
\]
\sub{Свободное падение из состояния покоя}
Пусть тело отпущено с высоты $h_0=9\unit{м}$ на неизвестной планете и упало за 3 с. Начало координат на поверхности, ось направлена вверх.
\[
 0=h_0-\frac{gt^2}{2}\quad\Rightarrow\quad
 g=\frac{2h_0}{t^2}=\frac{18}{9}=2\unit{м/с$^2$}.
\]
Если направите ось вниз, согласованно измените начальную координату и знак проекции ускорения.
\clearpage

\topic{4. Движение по окружности и невесомость}
Линейная скорость направлена по касательной, центростремительное ускорение --- к центру окружности. При радиусе $R>0$:
\[
 a_{\text{ц}}=\frac{v^2}{R},\qquad v=\sqrt{a_{\text{ц}}R},\qquad v=\omega R.
\]
$\omega$ --- угловая скорость (рад/с), $v$ --- модуль линейной скорости (м/с). Даже при постоянном модуле скорости её направление меняется.
\sub{Пример 1. Линейная скорость}
Если $R=50\unit{м}$ и $a_{\text{ц}}=2\unit{м/с$^2$}$, то
\[
 v=\sqrt{2\cdot50}=10\unit{м/с}.
\]
\sub{Пример 2. Мёртвая петля}
В верхней точке петли пилот испытывает невесомость. Вес --- сила, с которой тело действует на опору или подвес; в рассматриваемой задаче $P=0$ и $N=0$.
\begin{center}
\begin{tikzpicture}[scale=.82]
\draw[gray,thick] (0,0) circle (1.65);
\fill[accent] (0,1.65) circle (3pt);
\fill[gray] (0,0) circle (1.5pt) node[below left] {$O$};
\draw[dashed] (0,0)--(0,1.65);
\draw[vec] (0,1.65)--(1.45,1.65) node[right] {$\vec v$};
\draw[vec] (0,1.57)--(0,.5) node[left] {$m\vec g$};
\draw[vec] (.6,1.35)--(.6,.28) node[right] {$\vec a_{\text{ц}}$};
\draw[axis] (3,1.65)--(3,-.4) node[below] {$y$};
\node at (-2.6,1.4) {$N=0$};
\end{tikzpicture}
\end{center}
Направим ось вниз, к центру. В случае невесомости сила тяжести обеспечивает требуемое центростремительное ускорение:
\[
 mg+N=m\frac{v^2}{R},\qquad N=0
 \quad\Rightarrow\quad R=\frac{v^2}{g}.
\]
При $v=100\unit{м/с}$ и $g=10\unit{м/с$^2$}$ получаем $R=1000\unit{м}$.
\textbf{Важно:} невесомость означает отсутствие веса. Сила тяжести при этом может оставаться ненулевой. Центростремительное ускорение является результатом действия сил.
\clearpage

\topic{5. Сила упругости и соединение пружин}
Закон Гука для модуля силы упругости при линейной деформации:
\[
 F_{\text{упр}}=k|\Delta l|.
\]
$k$ --- жёсткость пружины (Н/м), $\Delta l$ --- изменение её длины (м). Сила направлена к восстановлению исходной длины.
\begin{center}
\begin{tikzpicture}[scale=.95]
\draw[thick] (-1.2,3.5)--(1.2,3.5);
\draw[spring] (0,3.4)--(0,2.1); \draw[thick] (0,2.1)--(0,1.9);
\draw[spring] (0,1.9)--(0,.6); \draw[thick] (0,.6)--(0,.35);
\draw[vec] (0,.35)--(0,-.25) node[right] {$\vec F$};
\node[right] at (.3,2.7) {$k_1$}; \node[right] at (.3,1.25) {$k_2$};
\node[below] at (0,-.48) {Последовательно};
\draw[thick] (4,3.5)--(6.6,3.5); \draw[spring] (4.6,3.4)--(4.6,.6); \draw[spring] (6,3.4)--(6,.6);
\draw[thick] (4.6,.6)--(6,.6); \draw[thick] (5.3,.6)--(5.3,.35);
\draw[vec] (5.3,.35)--(5.3,-.25) node[right] {$\vec F$};
\node[left] at (4.35,2) {$k_1$}; \node[right] at (6.25,2) {$k_2$};
\node[below] at (5.3,-.48) {Параллельно};
\end{tikzpicture}
\end{center}
\sub{Последовательное соединение}
Через обе идеальные пружины передаётся одна и та же сила; удлинения складываются. Поэтому
\[
 \frac1{k}=\frac1{k_1}+\frac1{k_2},\qquad k=\frac{k_1k_2}{k_1+k_2}.
\]
Пример: $k_1=60\unit{Н/м}$, $k_2=30\unit{Н/м}$.
\[
 \frac1k=\frac1{60}+\frac1{30}=\frac3{60}=\frac1{20},\qquad k=20\unit{Н/м}.
\]
\sub{Параллельное соединение}
При одинаковом удлинении пружин их силы складываются:
\[
 k=k_1+k_2=60+30=90\unit{Н/м}.
\]
\textbf{Самопроверка:} при $k_1,k_2>0$ последовательная система мягче каждой пружины, параллельная --- жёстче каждой.
\clearpage

\topic{6. Баланс сил и равномерное движение}
Второй закон Ньютона в векторной форме: $\sum\vec F=m\vec a$. При равномерном прямолинейном движении $\vec a=\vec0$, поэтому суммы проекций сил на обе оси равны нулю.
\sub{Пример 1. Трение при движении по горизонтали}
Коробку массой 10 кг тянут равномерно горизонтальной пружиной жёсткостью 200 Н/м. Удлинение пружины 0,2 м. Примем $g=10\unit{м/с$^2$}$.
\begin{center}
\begin{tikzpicture}[scale=.9]
\draw[thick] (-3,-.5)--(4.3,-.5);
\foreach \x in {-2.8,-2.4,...,4.2} \draw[gray] (\x,-.5)--++(-.2,-.2);
\draw[thick,fill=gray!8] (-.7,-.5) rectangle (.7,.5);

\fill[accent] (0,0) circle (2pt);
\draw[vec] (0,0)--(1.9,0) node[above] {$\vec F_{\text{упр}}$};
\draw[vec] (0,0)--(-1.9,0) node[above] {$\vec F_{\text{тр}}$};
\draw[vec] (0,0)--(0,1.5) node[left] {$\vec N$};
\draw[vec] (0,0)--(0,-1.5) node[left] {$m\vec g$};
\draw[axis] (3.4,.75)--(4.5,.75) node[right] {$x$};
\draw[axis] (3.4,.75)--(3.4,1.8) node[above] {$y$};
\end{tikzpicture}
\end{center}
Поскольку $a_x=a_y=0$, а пружина горизонтальна:
\[
 Ox:\ F_{\text{упр}}-F_{\text{тр}}=0,\qquad Oy:\ N-mg=0.
\]
Используем закон Гука и формулу трения скольжения $F_{\text{тр}}=\mu N$:
\[
 k\Delta l=\mu mg\quad\Rightarrow\quad
 \mu=\frac{k\Delta l}{mg}=\frac{200\cdot0{,}2}{10\cdot10}=0{,}4.
\]
Коэффициент трения безразмерный. Ускорение свободного падения $g$ входит в силу тяжести; ускорение самой коробки равно нулю.
\sub{Пример 2. Натяжение верёвки}
Девочка может тянуть с силой до 50 Н, мальчик --- до 150 Н. Для невесомой верёвки в равновесии силы на концах одинаковы по модулю. Максимальное натяжение ограничено меньшей доступной силой: $T=50\unit{Н}$.
\clearpage

\topic{Алгоритм: решение задачи по механике}
\vspace*{4mm}
\begin{center}
\begin{tikzpicture}[node distance=9mm]
\node[terminal] (start) {Начало};
\node[action,below=of start] (data) {Запишите «Дано», искомую величину; переведите данные в СИ.};
\node[action,below=of data] (draw) {Сделайте рисунок: тело, траектория, векторы; выберите оси и начало координат.};
\node[draw=accent,diamond,aspect=2.7,align=center,inner sep=1pt,text width=4.8cm,below=of draw] (choice) {Нужно учитывать силы?};
\node[action,text width=5.4cm,below left=14mm and -1.0cm of choice] (dyn) {Динамика: отметьте силы, запишите $\sum\vec F=m\vec a$ и проекции.};
\node[action,text width=5.4cm,below right=14mm and -1.0cm of choice] (kin) {Кинематика: выберите уравнения координаты, скорости или перемещения.};
\node[action,below=46mm of choice] (eq) {Свяжите уравнения с условиями задачи; выразите искомую величину.};
\node[action,below=of eq] (calc) {Подставьте числа, вычислите результат; проверьте знак, размерность и физический смысл.};
\node[terminal,below=of calc] (end) {Запишите ответ};
\draw[flowarrow] (start)--(data); \draw[flowarrow] (data)--(draw); \draw[flowarrow] (draw)--(choice);
\draw[flowarrow] (choice.west) -| node[pos=.25,above]{Да} (dyn.north);
\draw[flowarrow] (choice.east) -| node[pos=.25,above]{Нет} (kin.north);
\coordinate (merge) at ([yshift=5mm]eq.north);
\draw[draw=darkaccent,thick] (dyn.south) |- (merge); \draw[draw=darkaccent,thick] (kin.south) |- (merge);
\draw[flowarrow] (merge)--(eq.north);
\draw[flowarrow] (eq)--(calc); \draw[flowarrow] (calc)--(end);
\end{tikzpicture}
\end{center}
\clearpage

\topic{Тренировочный блок}
Выполняйте задачи с записью «Дано», рисунком и необходимыми уравнениями в проекциях. Если не оговорено иное, $g=10\unit{м/с$^2$}$, сопротивлением воздуха пренебрегайте.
\sub{Базовый уровень}
\begin{enumerate}[label=\arabic*.]
\item Шарик начал равноускоренное движение из покоя. За первые 2 с он прошёл 16 см. Найдите модуль ускорения.
\item Координата тела в начале движения равна $-12\unit{м}$, через 3 с она равна 8 м, через 6 с снова равна $-12\unit{м}$. Найдите модуль перемещения за 6 с.
\item Координата тела меняется по закону $x=-4+2t+3t^2$ (СИ). Найдите начальную координату, проекцию начальной скорости, проекцию ускорения и проекцию скорости в момент $t=4\unit{с}$.
\item Тело движется по окружности радиусом 32 м с центростремительным ускорением 8 м/с$^2$. Найдите модуль линейной скорости.
\item Две пружины имеют жёсткости 40 Н/м и 60 Н/м. Найдите жёсткость системы при последовательном и при параллельном соединении.
\end{enumerate}
\sub{Средний уровень}
\begin{enumerate}[label=\arabic*.,start=6]
\item Автомобиль движется прямолинейно и равномерно тормозит от 18 м/с до полной остановки за 12 с. Определите тормозной путь.
\item Мяч бросили вертикально вверх с поверхности земли со скоростью 12 м/с. Найдите время подъёма и максимальную высоту.
\item На неизвестной планете тело отпустили без начальной скорости с высоты 18 м. Оно достигло поверхности через 3 с. Найдите модуль ускорения свободного падения.
\end{enumerate}
\sub{Повышенный уровень}
\begin{enumerate}[label=\arabic*.,start=9]
\item В верхней точке мёртвой петли пилот испытывает невесомость при скорости самолёта 80 м/с. Найдите радиус петли.
\item Ящик массой 12 кг равномерно тянут по горизонтальной поверхности горизонтальной пружиной жёсткостью 300 Н/м. Пружина удлинена на 0,12 м. Найдите коэффициент трения скольжения.
\end{enumerate}
\clearpage

\topic{Ответы}
\begin{center}
{\small Таблица 1. Ответы к тренировочному блоку}\par\vspace{5pt}
\renewcommand{\arraystretch}{1.45}
\begin{tabularx}{.88\linewidth}{@{}lX@{}}
\toprule
№ & Ответ\\
\midrule
1 & $0{,}08\unit{м/с$^2$}$\\
2 & $0\unit{м}$\\
3 & $x_0=-4\unit{м}$; $v_{0x}=2\unit{м/с}$; $a_x=6\unit{м/с$^2$}$; $v_x(4)=26\unit{м/с}$\\
4 & $16\unit{м/с}$\\
5 & Последовательно: $24\unit{Н/м}$; параллельно: $100\unit{Н/м}$\\
6 & $108\unit{м}$\\
7 & $1{,}2\unit{с}$; $7{,}2\unit{м}$\\
8 & $4\unit{м/с$^2$}$\\
9 & $640\unit{м}$\\
10 & $0{,}3$\\
\bottomrule
\end{tabularx}
\end{center}
\vspace{8mm}
\topic{Итоговая памятка}
\begin{tabularx}{\linewidth}{@{}p{.38\linewidth}X@{}}
\toprule
Сигнал в задаче & Полезное действие\\
\midrule
Дано уравнение $x(t)$ & Сравните коэффициенты с общим уравнением координаты.\\
Нужно перемещение & Вычтите начальную координату из конечной; для модуля возьмите абсолютное значение.\\
Верхняя точка броска & Используйте $v_y=0$ и $a_y=-g$ при оси вверх.\\
Движение по окружности & Запишите $a_{\text{ц}}=v^2/R$; направление ускорения --- к центру.\\
Невесомость & Учтите $N=0$ в задаче о петле.\\
Равномерное движение & Запишите баланс сил: суммы проекций равны нулю.\\
Пружины & Определите тип соединения до вычисления общей жёсткости.\\
\bottomrule
\end{tabularx}
\end{document}
